% CODEX REPORT::CACHE-FINGERPRINT1::2026-10-07
\documentclass[11pt]{article}
\usepackage[margin=0.82in]{geometry}
\usepackage{amsmath,booktabs,graphicx,hyperref}
\title{CACHE-FINGERPRINT1: Exact Short and Adaptive Fingerprints in the Source-Replay WRAPPER}
\author{P versus NP Research Project}
\date{7 October 2026}

\begin{document}
\maketitle

\paragraph{Status.}
The cache-safety invariant below is \textbf{PROVED/EXACT}.  The measurements
are \textbf{COMPUTATIONAL}.  They establish neither an asymptotic cache bound
nor a general SAT contraction.  The guarded experiment completed 12 answer
checks with zero errors in 2.90 seconds, inside its 25-second internal limit.

\section{Mechanism and exactness guard}
The experiment uses the COMBINED-REPLAY1 solver unchanged outside its cache.
A state contains a source-clause scope and a partial assignment.  Residual
clauses are regenerated from the immutable source.  A fingerprint is computed
from the first-occurrence-normalized residual-clause stream.

For every candidate cache hit, the implementation regenerates both complete
normalized streams and compares them exactly before reusing the stored Boolean
answer.  Equal streams give a literal-sign-preserving variable bijection and
therefore equal satisfiability values.  An unequal stream is rejected.  Thus a
fingerprint collision can add work but cannot change an answer.  Collision
resistance of SHA-256 is not a correctness premise.

The fixed configurations retain 2, 4, 6, or 8 hexadecimal digits, respectively
8, 16, 24, or 32 bits.  The adaptive configurations start with either 2 or 4
hexadecimal digits and extend colliding entries by whole bytes, through at
most 8 hexadecimal digits.  Full 256-bit digests are transient.  They are
never hidden in cache entries.  When an old entry needs extension bits, its
residual is regenerated and rehashed, and that source work is charged.

Adaptive stored prefixes are separated whenever their digests differ within
the retained four-byte limit.  Prefixes still tied at that limit remain
candidates, whether the later digest bytes differ or the complete digest
collides, and the exact stream check continues to distinguish them.
Eviction can remove reuse opportunities only; it cannot introduce an answer.

\section{Charged experiment}
Each solver receives a 131,072-bit cache budget and uses frontier cap 5 and
block size 4.  The two workloads are the unsatisfiable PHP$(7,6)$ instance and
eight satisfiable, interleaved, isomorphic dense components of width five.
PHP supplies hundreds of distinct states and real short-prefix collisions;
the component instance checks successful equality reuse.  Every configuration
returns the known answer on both inputs.  All configurations make the same 892
recursive calls and obtain the same 121 verified cache hits.  No configuration
evicts an entry, so the table isolates fingerprint and verification costs.

The cache-bit column is the maximum cache descriptor payload in one solver.
It includes state descriptors, stored fingerprint prefixes, and answers.  The
workspace column is the maximum charged sequential-model peak, including live
recursion descriptors, cache, active wrapper data, normalization maps, and
hash temporaries.  Reads, candidate comparisons, and rejected collisions are
summed across the two workloads.

\begin{center}
\resizebox{\textwidth}{!}{%
\begin{tabular}{lrrrrrrrr}
\toprule
Fingerprint & Max cache & Bucket & Equality & Rejected & Literal & Recursive & Seq. peak\\
scheme & bits & peak & replays & collisions & reads & calls & bits\\
\midrule
Fixed 2 hex & 49,244 & 10 & 1,278 & 1,157 & 887,763 & 892 & 51,721\\
Fixed 4 hex & 55,396 & 2 & 122 & 1 & 833,920 & 892 & 57,577\\
Fixed 6 hex & 61,548 & 1 & 121 & 0 & 833,889 & 892 & 63,433\\
Fixed 8 hex & 67,700 & 1 & 121 & 0 & 833,889 & 892 & 69,289\\
Adaptive 2--8 hex & 55,116 & 10$^{*}$ & 312 & 191 & 989,748 & 892 & 57,273\\
Adaptive 4--8 hex & 55,412 & 2$^{*}$ & 122 & 1 & 834,328 & 892 & 57,593\\
\bottomrule
\end{tabular}}
\end{center}

For adaptive rows, the starred value is occupancy of the original short-prefix
bucket.  After extension, the maximum exact stored-fingerprint bucket occupancy
is one.  Random-access-model workspace peaks, which additionally charge the
source clause-boundary index, are respectively 52,918, 58,774, 64,630, 70,486,
58,470, and 58,790 bits in table order.

\begin{center}
\begin{tabular}{lrrrr}
\toprule
Adaptive scheme & Extended entries & Added prefix bits & Extension rehashes & Final prefix bits\\
\midrule
2--8 hex & 733 & 5,872 & 712 & 12,040\\
4--8 hex & 2 & 16 & 2 & 12,352\\
\bottomrule
\end{tabular}
\end{center}
The extension count includes both old entries extended in place and new entries
inserted with longer prefixes.  ``Extension rehashes'' counts regenerations of
old residuals needed to discover the extra bytes.

\section{Interpretation}
The 8-bit fixed cache minimizes stored bits in this batch, but its buckets reach
occupancy ten and it performs 1,157 rejected exact comparisons.  Moving to 16
bits eliminates all but one rejected collision while increasing the maximum
cache payload by 6,152 bits.  Moving from 16 to 24 bits removes that final
collision at the cost of another 6,152 cache bits; it saves only 31 literal
reads.  Moving from 24 to 32 bits changes neither equality work nor literal
reads and adds another 6,152 cache bits.  Fixed 16 bits is therefore the best
observed point for this bounded workload, not a universal optimum.

Starting adaptively at 8 bits cuts rejected collisions from 1,157 to 191 and
candidate equality replays from 1,278 to 312.  However, regeneration for 712
extension hashes raises literal reads by 11.5 percent relative to fixed 8 bits
and leaves the final prefix payload close to fixed 16 bits.  Starting at 16
bits extends only the two entries involved in the sole collision.  Since that
collision is not repeated later, adaptation cannot recover its extension cost
in this run.

The result supports a narrow engineering rule: exact short fingerprints are
safe when full residual replay remains mandatory, but their useful length is a
time--space choice.  Adaptive extension helps only when repeated collision
verification costs more than regenerating and retaining the extra prefix bits.
Nothing here supplies the missing uniform size reduction or total-branching
bound for the proof spine.

The executable artifact is \texttt{cache\_fingerprint\_experiment.py}; its
complete measurements are in \texttt{cache\_fingerprint\_results.json}.

% END_OF_VERIFIED_REPORT::CODEX::CACHE-FINGERPRINT1
\end{document}
